\documentclass[11pt,a4paper]{article}
\usepackage[a4paper,margin=24mm]{geometry}
\usepackage[utf8]{inputenc}
\usepackage[T1]{fontenc}
\usepackage{lmodern}
\usepackage{amsmath,amssymb,amsthm}
\usepackage{booktabs,array,tabularx}
\usepackage{enumitem}
\usepackage{float}
\usepackage{longtable}
\usepackage{microtype}
\usepackage{xcolor}
\definecolor{linkblue}{rgb}{0,0.15,0.5}
\usepackage[colorlinks=true,linkcolor=linkblue,citecolor=linkblue,urlcolor=linkblue]{hyperref}
\usepackage{url}
\allowdisplaybreaks
\emergencystretch=2em

\renewcommand{\thesection}{S\arabic{section}}
\renewcommand{\theequation}{S.\arabic{equation}}
\renewcommand{\thetable}{S\arabic{section}.\arabic{table}}
\theoremstyle{plain}
\newtheorem{theorem}{Theorem}[section]
\newtheorem{proposition}[theorem]{Proposition}
\newtheorem{lemma}[theorem]{Lemma}
\newtheorem{corollary}[theorem]{Corollary}
\theoremstyle{definition}
\newtheorem{definition}[theorem]{Definition}
\newtheorem{example}[theorem]{Example}
\theoremstyle{remark}
\newtheorem{remark}[theorem]{Remark}

\newcommand{\R}{\mathbb{R}}
\newcommand{\Qr}{\mathbb{Q}}
\newcommand{\Prob}{\mathbb{P}}
\newcommand{\calV}{\mathcal{V}}
\newcommand{\calS}{\mathcal{S}}
\newcommand{\calW}{\mathcal{W}}
\newcommand{\calD}{\mathcal{D}}
\newcommand{\Wfb}{\calW^{\mathrm{fb}}}
\newcommand{\Wbl}{\calW^{\mathrm{bl}}}
\newcommand{\supp}{\operatorname{supp}}
\newcommand{\Post}{\operatorname{Post}}
\newcommand{\hold}{\mathrm{hold}}
\newcommand{\ol}{\mathrm{ol}}
\newcommand{\seq}{\mathrm{seq}}
\newcommand{\obs}{\mathrm{obs}}
\newcommand{\Tobs}{T_{\mathrm{obs}}}
\newcommand{\Tl}{T_{\mathrm{learn}}}
\newcommand{\ind}[1]{\mathbf{1}[#1]}
\newcommand{\tfr}[2]{\tfrac{#1}{#2}}
\newcommand{\mref}[1]{\csname mainlab@#1\endcsname}
\newcommand{\meqref}[1]{(\mref{#1})}
\setlength{\parskip}{2pt}

\begin{document}

\title{\textbf{Supplementary material for}\\ \emph{Probabilistic sufficiency under partial observation: belief-state safety values, survivable-set antichains, and the exact price of late information}}
\author{Amin Abaee\\[0.3em]{\small Independent Researcher}\\[0.2em]{\small ORCID: \href{https://orcid.org/0000-0002-0019-1842}{0000-0002-0019-1842}\quad \href{mailto:amin_abaee@ut.ac.ir}{amin\_abaee@ut.ac.ir}}}
\date{October 15, 2026}
\maketitle

\noindent This supplement accompanies the main text (version 16 of the manuscript). Section, theorem and equation numbers without the prefix S refer to the main text. Every table below is generated by the script \texttt{make\_si\_tables.py} from the exact module \texttt{p02\_core.py}; no entry is typed by hand.

\tableofcontents

% generated from the main-text .aux file: main-text numbers used by \mref / \meqref
\expandafter\def\csname mainlab@sec:intro\endcsname{1}
\expandafter\def\csname mainlab@sec:setting\endcsname{2}
\expandafter\def\csname mainlab@def:model\endcsname{2.1}
\expandafter\def\csname mainlab@def:value\endcsname{2.2}
\expandafter\def\csname mainlab@prop:recursion\endcsname{2.3}
\expandafter\def\csname mainlab@eq:recursion\endcsname{1}
\expandafter\def\csname mainlab@prop:alpha\endcsname{2.4}
\expandafter\def\csname mainlab@eq:backup\endcsname{2}
\expandafter\def\csname mainlab@rem:complexity\endcsname{2.5}
\expandafter\def\csname mainlab@sec:support\endcsname{3}
\expandafter\def\csname mainlab@thm:support\endcsname{3.1}
\expandafter\def\csname mainlab@eq:support\endcsname{3}
\expandafter\def\csname mainlab@eq:minmass\endcsname{4}
\expandafter\def\csname mainlab@prop:feedback\endcsname{3.2}
\expandafter\def\csname mainlab@eq:voi\endcsname{5}
\expandafter\def\csname mainlab@thm:antichain\endcsname{3.3}
\expandafter\def\csname mainlab@rem:operators\endcsname{3.4}
\expandafter\def\csname mainlab@prop:freeze\endcsname{3.5}
\expandafter\def\csname mainlab@sec:deficit\endcsname{4}
\expandafter\def\csname mainlab@prop:deficit\endcsname{4.1}
\expandafter\def\csname mainlab@fig:deficits\endcsname{1}
\expandafter\def\csname mainlab@ex:twofloor\endcsname{4.2}
\expandafter\def\csname mainlab@ex:asym\endcsname{4.3}
\expandafter\def\csname mainlab@sec:delayed\endcsname{5}
\expandafter\def\csname mainlab@eq:delayed\endcsname{6}
\expandafter\def\csname mainlab@thm:closed\endcsname{5.1}
\expandafter\def\csname mainlab@eq:hold\endcsname{7}
\expandafter\def\csname mainlab@eq:ol\endcsname{8}
\expandafter\def\csname mainlab@eq:obs\endcsname{9}
\expandafter\def\csname mainlab@thm:lattice\endcsname{5.2}
\expandafter\def\csname mainlab@fig:closed\endcsname{2}
\expandafter\def\csname mainlab@eq:flex\endcsname{10}
\expandafter\def\csname mainlab@eq:voi2\endcsname{11}
\expandafter\def\csname mainlab@cor:timing\endcsname{5.3}
\expandafter\def\csname mainlab@rem:declaration\endcsname{5.4}
\expandafter\def\csname mainlab@cor:grid\endcsname{5.5}
\expandafter\def\csname mainlab@sec:deadline\endcsname{6}
\expandafter\def\csname mainlab@eq:declining\endcsname{12}
\expandafter\def\csname mainlab@fig:deadlines\endcsname{3}
\expandafter\def\csname mainlab@prop:deadline\endcsname{6.1}
\expandafter\def\csname mainlab@eq:deadline\endcsname{13}
\expandafter\def\csname mainlab@thm:additive\endcsname{6.2}
\expandafter\def\csname mainlab@eq:general\endcsname{14}
\expandafter\def\csname mainlab@eq:additive\endcsname{15}
\expandafter\def\csname mainlab@ex:additive\endcsname{6.3}
\expandafter\def\csname mainlab@rem:delayedschedule\endcsname{6.4}
\expandafter\def\csname mainlab@sec:stochastic\endcsname{7}
\expandafter\def\csname mainlab@prop:pathdeficit\endcsname{7.1}
\expandafter\def\csname mainlab@eq:pathdeficit\endcsname{16}
\expandafter\def\csname mainlab@eq:sensor\endcsname{17}
\expandafter\def\csname mainlab@prop:noisyprobe\endcsname{7.2}
\expandafter\def\csname mainlab@eq:noisyprobe\endcsname{18}
\expandafter\def\csname mainlab@eq:majority3\endcsname{19}
\expandafter\def\csname mainlab@fig:probes\endcsname{4}
\expandafter\def\csname mainlab@prop:probecount\endcsname{7.3}
\expandafter\def\csname mainlab@eq:probecount\endcsname{20}
\expandafter\def\csname mainlab@rem:eps0\endcsname{7.4}
\expandafter\def\csname mainlab@sec:dr\endcsname{8}
\expandafter\def\csname mainlab@eq:rho\endcsname{21}
\expandafter\def\csname mainlab@prop:dr\endcsname{8.1}
\expandafter\def\csname mainlab@eq:contam\endcsname{22}
\expandafter\def\csname mainlab@prop:priorcontam\endcsname{8.2}
\expandafter\def\csname mainlab@eq:priorcontam\endcsname{23}
\expandafter\def\csname mainlab@sec:related\endcsname{9}
\expandafter\def\csname mainlab@sec:verification\endcsname{10}
\expandafter\def\csname mainlab@sec:discussion\endcsname{11}
\section{The certified value table and its census}\label{si:grid}

Table~\ref{tab:grid} is the complete hold-class value table on the certified grid of the delayed hidden-regime instance (Section~\mref{sec:delayed} of the main text) at the symmetric prior: $z_0\in\{1.0,\dots,2.5\}$, $\Tobs\in\{1,2,3\}$, horizons $k\in\{1,\dots,4\}$, $192$ cell--horizon combinations. Each entry is computed twice, by segment-form enumeration of the two holds and by the closed form \meqref{eq:hold} of Theorem~\mref{thm:closed}, and the two agree on every combination.

\begin{table}[h]
\centering\footnotesize
\caption{Hold-class safety value $V^{\hold}_k(b_0)$ on the certified grid at $w=\tfr12$. The value is $1$ exactly when $z_0\ge1+\min(k,\Tobs)$: $36$ unit entries and $156$ half entries.}
\label{tab:grid}
\setlength{\tabcolsep}{4pt}
\begin{tabular}{@{}l cccc c cccc c cccc@{}}
\toprule
 & \multicolumn{4}{c}{$\Tobs=1$} & & \multicolumn{4}{c}{$\Tobs=2$} & & \multicolumn{4}{c}{$\Tobs=3$}\\
\cmidrule{2-5}\cmidrule{7-10}\cmidrule{12-15}
$z_0$ & $k{=}1$ & $k{=}2$ & $k{=}3$ & $k{=}4$ & & $k{=}1$ & $k{=}2$ & $k{=}3$ & $k{=}4$ & & $k{=}1$ & $k{=}2$ & $k{=}3$ & $k{=}4$\\
\midrule
1.0 & $\tfrac{1}{2}$ & $\tfrac{1}{2}$ & $\tfrac{1}{2}$ & $\tfrac{1}{2}$ & & $\tfrac{1}{2}$ & $\tfrac{1}{2}$ & $\tfrac{1}{2}$ & $\tfrac{1}{2}$ & & $\tfrac{1}{2}$ & $\tfrac{1}{2}$ & $\tfrac{1}{2}$ & $\tfrac{1}{2}$\\
1.1 & $\tfrac{1}{2}$ & $\tfrac{1}{2}$ & $\tfrac{1}{2}$ & $\tfrac{1}{2}$ & & $\tfrac{1}{2}$ & $\tfrac{1}{2}$ & $\tfrac{1}{2}$ & $\tfrac{1}{2}$ & & $\tfrac{1}{2}$ & $\tfrac{1}{2}$ & $\tfrac{1}{2}$ & $\tfrac{1}{2}$\\
1.2 & $\tfrac{1}{2}$ & $\tfrac{1}{2}$ & $\tfrac{1}{2}$ & $\tfrac{1}{2}$ & & $\tfrac{1}{2}$ & $\tfrac{1}{2}$ & $\tfrac{1}{2}$ & $\tfrac{1}{2}$ & & $\tfrac{1}{2}$ & $\tfrac{1}{2}$ & $\tfrac{1}{2}$ & $\tfrac{1}{2}$\\
1.3 & $\tfrac{1}{2}$ & $\tfrac{1}{2}$ & $\tfrac{1}{2}$ & $\tfrac{1}{2}$ & & $\tfrac{1}{2}$ & $\tfrac{1}{2}$ & $\tfrac{1}{2}$ & $\tfrac{1}{2}$ & & $\tfrac{1}{2}$ & $\tfrac{1}{2}$ & $\tfrac{1}{2}$ & $\tfrac{1}{2}$\\
1.4 & $\tfrac{1}{2}$ & $\tfrac{1}{2}$ & $\tfrac{1}{2}$ & $\tfrac{1}{2}$ & & $\tfrac{1}{2}$ & $\tfrac{1}{2}$ & $\tfrac{1}{2}$ & $\tfrac{1}{2}$ & & $\tfrac{1}{2}$ & $\tfrac{1}{2}$ & $\tfrac{1}{2}$ & $\tfrac{1}{2}$\\
1.5 & $\tfrac{1}{2}$ & $\tfrac{1}{2}$ & $\tfrac{1}{2}$ & $\tfrac{1}{2}$ & & $\tfrac{1}{2}$ & $\tfrac{1}{2}$ & $\tfrac{1}{2}$ & $\tfrac{1}{2}$ & & $\tfrac{1}{2}$ & $\tfrac{1}{2}$ & $\tfrac{1}{2}$ & $\tfrac{1}{2}$\\
1.6 & $\tfrac{1}{2}$ & $\tfrac{1}{2}$ & $\tfrac{1}{2}$ & $\tfrac{1}{2}$ & & $\tfrac{1}{2}$ & $\tfrac{1}{2}$ & $\tfrac{1}{2}$ & $\tfrac{1}{2}$ & & $\tfrac{1}{2}$ & $\tfrac{1}{2}$ & $\tfrac{1}{2}$ & $\tfrac{1}{2}$\\
1.7 & $\tfrac{1}{2}$ & $\tfrac{1}{2}$ & $\tfrac{1}{2}$ & $\tfrac{1}{2}$ & & $\tfrac{1}{2}$ & $\tfrac{1}{2}$ & $\tfrac{1}{2}$ & $\tfrac{1}{2}$ & & $\tfrac{1}{2}$ & $\tfrac{1}{2}$ & $\tfrac{1}{2}$ & $\tfrac{1}{2}$\\
1.8 & $\tfrac{1}{2}$ & $\tfrac{1}{2}$ & $\tfrac{1}{2}$ & $\tfrac{1}{2}$ & & $\tfrac{1}{2}$ & $\tfrac{1}{2}$ & $\tfrac{1}{2}$ & $\tfrac{1}{2}$ & & $\tfrac{1}{2}$ & $\tfrac{1}{2}$ & $\tfrac{1}{2}$ & $\tfrac{1}{2}$\\
1.9 & $\tfrac{1}{2}$ & $\tfrac{1}{2}$ & $\tfrac{1}{2}$ & $\tfrac{1}{2}$ & & $\tfrac{1}{2}$ & $\tfrac{1}{2}$ & $\tfrac{1}{2}$ & $\tfrac{1}{2}$ & & $\tfrac{1}{2}$ & $\tfrac{1}{2}$ & $\tfrac{1}{2}$ & $\tfrac{1}{2}$\\
2.0 & $1$ & $1$ & $1$ & $1$ & & $1$ & $\tfrac{1}{2}$ & $\tfrac{1}{2}$ & $\tfrac{1}{2}$ & & $1$ & $\tfrac{1}{2}$ & $\tfrac{1}{2}$ & $\tfrac{1}{2}$\\
2.1 & $1$ & $1$ & $1$ & $1$ & & $1$ & $\tfrac{1}{2}$ & $\tfrac{1}{2}$ & $\tfrac{1}{2}$ & & $1$ & $\tfrac{1}{2}$ & $\tfrac{1}{2}$ & $\tfrac{1}{2}$\\
2.2 & $1$ & $1$ & $1$ & $1$ & & $1$ & $\tfrac{1}{2}$ & $\tfrac{1}{2}$ & $\tfrac{1}{2}$ & & $1$ & $\tfrac{1}{2}$ & $\tfrac{1}{2}$ & $\tfrac{1}{2}$\\
2.3 & $1$ & $1$ & $1$ & $1$ & & $1$ & $\tfrac{1}{2}$ & $\tfrac{1}{2}$ & $\tfrac{1}{2}$ & & $1$ & $\tfrac{1}{2}$ & $\tfrac{1}{2}$ & $\tfrac{1}{2}$\\
2.4 & $1$ & $1$ & $1$ & $1$ & & $1$ & $\tfrac{1}{2}$ & $\tfrac{1}{2}$ & $\tfrac{1}{2}$ & & $1$ & $\tfrac{1}{2}$ & $\tfrac{1}{2}$ & $\tfrac{1}{2}$\\
2.5 & $1$ & $1$ & $1$ & $1$ & & $1$ & $\tfrac{1}{2}$ & $\tfrac{1}{2}$ & $\tfrac{1}{2}$ & & $1$ & $\tfrac{1}{2}$ & $\tfrac{1}{2}$ & $\tfrac{1}{2}$\\
\bottomrule
\end{tabular}
\end{table}

\paragraph{Census.} The $36$ unit entries are the six states $z_0\ge2.0$ at all four horizons for $\Tobs=1$ ($24$) and at $k=1$ for $\Tobs=2,3$ ($12$). Each combination carries two hold-direction witnesses, $\bigl(\ind{\text{branch }-\text{ survives}},\ind{\text{branch }+\text{ survives}}\bigr)$ for $u\equiv+1$ and $u\equiv-1$: on unit-value combinations both equal $(1,1)$ ($72$ occurrences), on half-value combinations they are the mirror pair $(1,0)$ and $(0,1)$ ($156$ each). The total is $384$ vectors; exact-equality deduplication within each combination leaves $348$ ($36+2\times156$). The three types are the indicators of the maximal survivable sets of Theorem~\mref{thm:antichain}: $\{+,-\}$ when $z_0\ge1+\ell$ and the two singletons otherwise.

\paragraph{Partition.} The certificate partition of the $48$ cells coincides with the value boundaries: $6$ viable cells ($\Tobs=1$, $z_0\ge2.0$: $V_k\equiv1$), $30$ common-action cells ($z_0<2.0$: $V_1<1$) and $12$ timing cells ($\Tobs\in\{2,3\}$, $z_0\ge2.0$: $V_1=1$ but $V_{\Tobs+1}<1$).

\paragraph{Class difference.} The open-loop value \meqref{eq:ol} equals the hold value except on $\{2\le z_0<1+\min(k,\Tobs)\}$; on the grid this is the $12$ timing cells at each of $k=2,3,4$, i.e.\ $36$ combinations. The identity of the brute-force difference set with the region is verified on the extended grid $z_0\in\{1.0,\dots,4.0\}$, where the region is no longer a union of whole timing-cell columns (for instance at $z_0=3.5$, $\Tobs=3$, $k=2$ both classes attain one).
\section{The two-floor instance and the asymmetric audit}\label{si:instances}

\paragraph{Two-floor instance.} States $\{0,1\}$ and actions $u\in\{\tfr25,\tfr35\}$; state $0$ stays at $0$ under $u=\tfr25$ and exits under $u=\tfr35$, state $1$ stays at $1$ under $u=\tfr35$ and exits under $u=\tfr25$; a single uninformative observation. The masked backup gives $\Gamma_1=\{(1,0),(0,1)\}$ and, since each surviving state is kept alive indefinitely by its own action, $\Gamma_k=\Gamma_1$ and $V_k(b)=\max(b_0,b_1)$ for every $k\ge1$ (verified to $k=5$). The one-step exit probabilities are $p(0,\tfr35)=p(1,\tfr25)=1$, so at $b=(\tfr12,\tfr12)$ the chance-constrained certificate $\sum_xb(x)p(x,a)\ge\tfr12$ fires for both actions, giving $\Delta_k\ge\tfr12$, which is attained. The witness vectors live on $\calV\cup\{\bot\}$ as $(1,0,0)$ and $(0,1,0)$; the suppressed bottom coordinate is identically zero.

\paragraph{Asymmetric audit.} On the delayed instance at $(z_0,\Tobs)=(\tfr32,2)$ with prior mass $w=\tfr3{10}$ on $\theta=+1$, Table~\ref{tab:asym} lists the four open-loop two-step segments. The first move sends one branch to $\tfr12<1$, so every segment loses a branch at step $1$; the best segments save the heavier branch and $V_2(b_0)=\tfr7{10}$, deficit $\tfr3{10}=\min(w,1-w)$.

\begin{table}[h]
\centering\footnotesize
\caption{The asymmetric audit at $(z_0,\Tobs)=(\tfr32,2)$, $b_0=(w,1-w)=(\tfr3{10},\tfr7{10})$.}
\label{tab:asym}
\begin{tabular}{@{}l c c l@{}}
\toprule
Segment $(u_0,u_1)$ & positions after step 1 ($\theta=+1$, $\theta=-1$) & surviving mass & lost branch\\
\midrule
$(+1,+1)$ & $5/2$, $1/2$ & $\tfrac{3}{10}$ & $\theta=-1$ (mass $\tfrac{7}{10}$), at step 1\\
$(+1,-1)$ & $5/2$, $1/2$ & $\tfrac{3}{10}$ & $\theta=-1$ (mass $\tfrac{7}{10}$), at step 1\\
$(-1,+1)$ & $1/2$, $5/2$ & $\tfrac{7}{10}$ & $\theta=+1$ (mass $\tfrac{3}{10}$), at step 1\\
$(-1,-1)$ & $1/2$, $5/2$ & $\tfrac{7}{10}$ & $\theta=+1$ (mass $\tfrac{3}{10}$), at step 1\\
\bottomrule
\end{tabular}
\end{table}
\section{Jump loci and agreement of the segment enumeration}\label{si:loci}

\paragraph{Agreement.} The declared hold class couples the stages of the blind window, so its value is not computed by the stagewise recursion over $A=\{\pm1\}$---that recursion computes the sequential value, which equals the open-loop value \meqref{eq:ol} and is one on the timing cells. The segment form of Proposition~\mref{prop:alpha} enumerates the two admissible holds and evaluates the surviving mass of each; Table~\ref{tab:grid} records that this enumeration returns \meqref{eq:hold} on all $192$ combinations. The class constraint is load-bearing: Corollary~\mref{cor:timing} is a statement about which class the timing certificate speaks for.

\paragraph{Jump loci.} Fix $\Tobs$, $k\ge1$ and $w$, and read \meqref{eq:hold} as a function of the real parameter $z_0$. It is piecewise constant with exactly two jumps: at $z_0=1$, from $0$ to $\max(w,1-w)$, and at $z_0=1+\ell$, $\ell=\min(k,\Tobs)$, from $\max(w,1-w)$ to $1$; at the symmetric prior both jumps have size $\tfr12$. For $k\le\Tobs$ the second locus is $z_0=1+k$---an off-grid consequence of the closed form (for $k=2$ and $\Tobs\ge2$ the jump is at $z_0=3$, not at $2$), verified on the extended grid. The witness family changes at the same loci: at the one-step horizon both mirror witnesses $(1,0)$ and $(0,1)$ attain the value $\max(w,1-w)$ below $z_0=2$ and the merged witness $(1,1)$ attains $1$ at and above it; at horizon $k$ the change sits at $z_0=1+\ell$. These are jumps in a parameter of the instance; as a function of the belief $b$ the value is piecewise linear and convex (Proposition~\mref{prop:alpha}), its kinks being the usual ones of the $\alpha$-vector representation.

\paragraph{Deficit profile.} On the $42$ nonviable cells the hold-class deficit at horizon $k$ is $\tfr12$ if $k>z_0-1$ and $0$ otherwise (for $k\le\Tobs$), and $\tfr12$ for all $k>\Tobs$: on the timing cells $0$ at $k=1$ and $\tfr12$ at $k\ge2$, on the common-action cells $\tfr12$ at every $k\ge1$. The deficit is non-decreasing in $k$ (Proposition~\mref{prop:deficit}(iii)) and frozen from $k=\Tobs$ on (Proposition~\mref{prop:freeze}); the open-loop deficit is frozen from $k=1$. Both freezes are verified by enumeration to horizon $40$.
\section{Learning deadlines and the two-patch instance}\label{si:deadline}

\paragraph{Deadline table.} Table~\ref{tab:deadline} lists, for each deadline $\Tl=0,\dots,4$, the exact pinned-probe threshold $\tfr{21}{10}+\tfr{\Tl}{10}$ of Proposition~\mref{prop:deadline}, the resulting viable set on the grid, and the probe-free benchmark $1+\tfr{\Tl}{10}$. All thresholds are verified by simulating both branches to horizon $40$ for every grid state; the free-probe threshold is $\tfr{21}{10}$ at every deadline. The second aligned instance of Example~\mref{ex:additive}, $z_{t+1}=z_t-\tfr3{20}+\tfr9{10}\theta u_t$, has thresholds $\tfr{41}{20}+\tfr{3T}{20}$, verified on a grid of step $\tfr1{20}$ for $T=0,\dots,6$; the misaligned instance has exact thresholds $1+\max\{T,\,1\}$ for $T\ge0$ against the additive value $1+T+1$.

\begin{table}[h]
\centering\footnotesize
\caption{Pinned-probe learning deadlines on the declining instance and the probe-free benchmark.}
\label{tab:deadline}
\begin{tabular}{@{}c c c l c c@{}}
\toprule
$\Tl$ & pinned threshold & $|K|$ & kernel $K$ on the grid & probe-free threshold & $|K|$ (probe-free)\\
\midrule
0 & $2.1$ & 5 & $\{2.1, 2.2, 2.3, 2.4, 2.5\}$ & $1.0$ & 16\\
1 & $2.2$ & 4 & $\{2.2, 2.3, 2.4, 2.5\}$ & $1.1$ & 15\\
2 & $2.3$ & 3 & $\{2.3, 2.4, 2.5\}$ & $1.2$ & 14\\
3 & $2.4$ & 2 & $\{2.4, 2.5\}$ & $1.3$ & 13\\
4 & $2.5$ & 1 & $\{2.5\}$ & $1.4$ & 12\\
\bottomrule
\end{tabular}
\end{table}

\paragraph{Two-patch instance with unknown regeneration.} Two patches carry integer stocks $(x_1,x_2)\in\{0,\dots,3\}^2$, capped at $3$, with a common unknown regeneration rate $g\in\{1,2\}$; in each stage the manager spares one patch $u\in\{1,2\}$, which grows by $g$, while the other grows by $g$ and is harvested by $2$. The safe set is $\{x_1\ge1,x_2\ge1\}$. Consider the belief with support $\{(1,2),(2,1)\}\times\{1,2\}$ and an observation that reports only the aggregate stock $x_1+x_2$. The feedback recursion to depth $12$ (and $10$, for stability) gives: nonviable when $g$ is hidden and only the aggregate is observed; viable when $(x,g)$ is observed per cell; and viable for each fixed $g$ under full state observation. The instance illustrates the qualitative statement behind Section~\mref{sec:deadline}: a hidden regeneration rate is a structural, not a measurement, obstruction---no amount of aggregate monitoring removes it, and revelation of $g$ does.
\section{Exact wrong-majority probabilities}\label{si:majority}

For $n=2m+1$ independent readings with error $\varepsilon$ the wrong-majority probability is the binomial tail $p_{\mathrm{wrong}}(n)=\sum_{j=m+1}^{n}\binom nj\varepsilon^{j}(1-\varepsilon)^{n-j}$. For $n=3$ this is $3\varepsilon^2(1-\varepsilon)+\varepsilon^3=\varepsilon^2(3-2\varepsilon)$, equation \meqref{eq:majority3}; at $\varepsilon=\tfr1{10}$ it equals $\tfr{7}{250}$. Table~\ref{tab:majority} lists the exact values at $\varepsilon=\tfr1{10}$ against the bound $(4\varepsilon(1-\varepsilon))^{n/2}=(3/5)^n$ of Proposition~\mref{prop:probecount}; the bound is loose by a factor between $6$ and $13$ on this range but has the correct exponential rate up to the constant in the exponent ($\ln(5/3)\approx0.51$ per reading against the empirical $\approx0.62$).

\begin{table}[h]
\centering\footnotesize
\caption{Exact wrong-majority probabilities at $\varepsilon=\tfr1{10}$ and the bound $(3/5)^n$.}
\label{tab:majority}
\begin{tabular}{@{}c c c c@{}}
\toprule
$n$ & exact $p_{\mathrm{wrong}}(n)$ & decimal & bound $(3/5)^n$\\
\midrule
1 & $\tfrac{1}{10}$ & 0.1 & 0.6\\
3 & $\tfrac{7}{250}$ & 0.028 & 0.216\\
5 & $\tfrac{107}{12500}$ & 0.00856 & 0.07776\\
7 & $\tfrac{341}{125000}$ & 0.002728 & 0.02799\\
9 & $\tfrac{22273}{25000000}$ & 0.00089092 & 0.01008\\
11 & $\tfrac{1848163}{6250000000}$ & 0.000295706 & 0.003628\\
\bottomrule
\end{tabular}
\end{table}

\paragraph{Majority-of-three value.} With readings at $t_p,t_p+1,t_p+2$ and $u=0$ in between, the state at the decision stage is $z_0-\tfr{t_p+2}{10}$ and the value is $(1-p_3)\ind{z_0-\tfr{t_p+2}{10}\ge1}+p_3\ind{z_0-\tfr{t_p+2}{10}\ge1+\tfr{11}{10}(k-t_p-2)}$; the script re-derives it by enumerating the eight reading patterns on both branches for $z_0$ on a grid of step $\tfr1{10}$, $t_p\in\{0,\dots,3\}$ and $k\le t_p+8$.
\section{The contamination operator: conventions and computations}\label{si:contam}

\paragraph{Conventions.} The contaminated recursion \meqref{eq:rho} takes the infimum over $q=(1-\rho)p+\rho r$ with $r$ ranging over all distributions on the outcome set $Y^{\bot}$ of the stage. The adversary may place mass on an outcome $y$ that is nominally null, $\Prob(y\mid b,a)=0$, for which the Bayes posterior is undefined; the operator therefore needs an explicit continuation there. The exact module adopts the convention that a nominally null outcome in $Y$ carries continuation value $0$ (the most pessimistic reading) and that the exit outcome $y_\bot$ always carries value $0$; the main-text statements are insensitive to the choice because in the instances used every outcome in the contaminating support is nominally positive. Fixing the nominal posteriors while perturbing the outcome probabilities makes \meqref{eq:rho} a rectangular stress operator in the sense of robust dynamic programming; it is not the Bayesian update under one coherent perturbed joint law, and no such interpretation is claimed.

\paragraph{The one-stage identity.} Equation \meqref{eq:contam} is the minimum of the linear functional $r\mapsto(1-\rho)\sum_yp(y)v_y+\rho\sum_yr(y)v_y$ over the simplex, attained at a vertex. The module verifies it against explicit vertex enumeration on the instances of the main text.

\paragraph{The non-affine example.} One safe state $s$, one action, $T(s\mid s)=\tfr12$, observations $y_s$ (after a safe transition) and $y_\bot$ (after exit), both nominally of probability $\tfr12$ from $\delta_s$. Then $V^\rho_1(\delta_s)=(1-\rho)\bigl(\tfr12\cdot1+\tfr12\cdot0\bigr)+\rho\min\{1,0\}=\tfr{1-\rho}2$ and, inductively, $V^\rho_{k}(\delta_s)=\bigl(\tfr{1-\rho}2\bigr)^{k}$. At $\rho=\tfr12$ and $k=2$ this is $\tfr1{16}$, whereas the affine interpolation of the endpoint values $V^0_2=\tfr14$ and $V^1_2=0$ gives $\tfr18$. The module computes $V^{1/2}_2$ by the recursion itself.

\paragraph{Prior contamination on the delayed class.} In \meqref{eq:priorcontam} the adversary chooses $r\in[0,1]$ after the hold is announced. On the band $1\le z_0<1+\ell$ the hold $u\equiv+1$ saves the $\theta=+1$ branch only, of contaminated mass $(1-\rho)w+\rho r\ge(1-\rho)w$, and the hold $u\equiv-1$ saves the other branch, of mass $(1-\rho)(1-w)+\rho(1-r)\ge(1-\rho)(1-w)$; the controller's best guaranteed value is $(1-\rho)\max(w,1-w)$. The module verifies the closed form by enumerating the two holds against the two extreme contaminations for $\rho\in\{0,\tfr14,\tfr12,\tfr34,1\}$, $w\in\{0,\tfr1{10},\dots,1\}$ and the grid of $(z_0,\ell)$.

\paragraph{Convexity.} At $\rho=0$ the value is convex (Proposition~\mref{prop:alpha}); for $\rho>0$ the inner minimum is a concave operation and convexity is lost in general. Concavity never holds: the two-floor value $\max(b_0,b_1)$ at $\rho=0$ equals $1$ at both vertices and $\tfr12$ at the midpoint.
\section{Verification record}\label{si:verification}

\paragraph{Files.} The edition is reproduced by four scripts, all in exact rational arithmetic with the Python standard library (plus \texttt{matplotlib} for the figures only):
\begin{itemize}[leftmargin=1.6em,itemsep=1pt]
\item \texttt{p02\_core.py} --- the exact module: a generic finite safety POMDP (masked $\alpha$-vector backup \meqref{eq:backup} with dominance pruning, Bayes update, open-loop enumeration, contaminated operator \meqref{eq:rho}); the delayed instance (closed forms, brute-force enumeration of holds and segments, census, cell classes, class difference); the declining instance (pinned, free and probe-free deadlines, the general schedule law, noisy probes, majority reading, probe-count tail); the feedback family $\mathcal F_n$; the two-floor, sensor, contamination and two-patch instances.
\item \path{paper02_probabilistic_sufficiency_v16_verification.py} --- $62$ checks (listed below); exits with status $0$ only if all pass. Its output is archived as \path{paper02_probabilistic_sufficiency_v16_verification_output.txt}.
\item \texttt{make\_figs.py} --- the four figures, every plotted value recomputed from the module.
\item \texttt{make\_si\_tables.py} --- Tables~\ref{tab:grid}, \ref{tab:asym}, \ref{tab:deadline} and \ref{tab:majority}, with built-in assertions against the closed forms.
\end{itemize}
The earlier campaign scripts (\path{paper2_belief_state_v2_verification.py}, which chains \path{paper2_stochastic_selector_v2_verify.py} and \path{hidden_parameter_learning_v1_verify.py}, and the figure script \path{paper2_belief_state_figures.py}) remain in the repository under \path{arena agent 1/paper rewrites/latex}; their checks are superseded by, and consistent with, the list below.

\paragraph{Check list.} S2.1--S2.3 two-floor witnesses and values; S4.1 chance-constrained certificate attained; S3.1--S3.2 the family $\mathcal F_n$ ($n=2,3,4$, random priors) and the strict feedback/blind separation; S3.3 witnesses are $\{0,1\}$ antichains; S3.4 the Sperner bound attained at $n=4$; S5.1--S5.2 closed forms against brute force on $1386$ structured triples $(z_0,w,\ell)$ including $z_0<1$ and a seeded battery of $500$ random cases; S5.3 the generic $\alpha$-vector backup on the finite reachable model reproduces the open-loop closed form; S5.4 the observed value; S5.5 the class lattice and both gaps; S5.6--S5.8 census, witnesses and partition; S5.9--S5.11 the class-difference set on the certified and extended grids and the off-grid example; S5.12 jump loci; S5.13--S5.14 the deficit profile; S4.2 the asymmetric audit; S3.5--S3.6 stabilization to horizon $40$ and nestedness; S6.1--S6.5 the deadline thresholds, kernel counts and benchmarks on a $\tfr1{20}$ grid; S6.6--S6.8 the general schedule law on both aligned instances, the misaligned example and the delayed instance as a schedule; S7.1--S7.3 the sensor instance; S7.4--S7.6 the noisy-probe and majority closed forms against exhaustive enumeration and the three value levels; S7.7--S7.10 the majority probabilities, the exact tail values, the bound for four error rates and the middle-band deficit; S8.1--S8.5 the contamination example, the one-stage identity against $200$ random vertex enumerations, the prior-contamination closed form and convexity; SI4.1 the two-patch instance; TXT.1--TXT.4 the compiled theorem numbers quoted in figure titles and captions, read from the \texttt{.aux} file; the presence of the headline tokens in the manuscript source; and TXT.5 the agreement of this supplement's table of main-text numbers with the compiled auxiliary file.

\paragraph{Lean layer: provenance.} The machine-checked statements live in the repository's \texttt{P3} modules (Lean~4, project pin \texttt{v4.34.1}). The committed pin was rebuilt in GitHub Actions (default Lake target, $60/60$ jobs, exit $0$; run \texttt{36946730855}); the exact project source and two footprint-check files are archived with SHA-256 manifests. A comment-aware scan of all $62$ source files finds no \texttt{sorry}, \texttt{admit}, explicit \texttt{axiom} or \texttt{constant} declaration; twenty-one named footprints contain no \texttt{sorryAx}; the aggregate-import environment exposes three source-level \texttt{native\_decide} uses in the support/value examples, yielding nine distinct generated axiom names under this pin---executable computational certificates that extend kernel-only trust to the Lean compiler and runtime; they are not admitted gaps and are not described as zero axioms or kernel-only proofs. The pinned declared-exceptions gate passed (run \texttt{36990482599}): only those nine names from three approved sites are permitted, and any change to those sites requires separate review. Three preserved historical modules outside the aggregate import closure are not certified by the default build. The declarations relevant to this paper are: \texttt{support\_identity} and \texttt{repairA\_feedback} (the support identity \meqref{eq:support} and $V_k(b)=1\iff\supp b\in\Wfb_k$ under deterministic kernels, with \texttt{paper\_literal\_thm\_support\_fails} refuting the earlier blind reading by an explicit counterexample), \texttt{VFb\_min\_mass\_bound} (the bound \meqref{eq:minmass}), \texttt{Vfam\_prune\_eq} and \texttt{VfamAdm\_frozen\_eq\_pruned} (pruning and the frozen value of Proposition~\mref{prop:freeze} for the specified families), \texttt{VR\_eq\_max\_alpha}, \texttt{VAdm\_convex} and the rationality lemmas (Proposition~\mref{prop:alpha}), \texttt{class\_monotone} (the first assertion of Theorem~\mref{thm:lattice}), \texttt{path\_deficit} (the fixed-policy form of \meqref{eq:pathdeficit}) and \texttt{contam\_lower}/\texttt{contam\_sharp} (the one-stage identity \meqref{eq:contam} with a point-mass minimizer). Not formalized: the Sperner bound of Theorem~\mref{thm:antichain}(iii), the parametric closed forms as such, Theorem~\mref{thm:additive}, the noisy-probe formulas and the probe-count bound; these are verified by exact enumeration only.

\paragraph{Figure provenance.} Figure~\mref{fig:deficits}: \texttt{fig1\_deficits.pdf} (two-floor value and sensor deficits); Figure~\mref{fig:closed}: \texttt{fig2\_closedforms.pdf} (closed forms and the certified-grid class map); Figure~\mref{fig:deadlines}: \texttt{fig3\_deadlines.pdf}; Figure~\mref{fig:probes}: \texttt{fig4\_probes.pdf}. All are produced by \texttt{make\_figs.py} from \texttt{p02\_core.py}; the result numbers in their titles are checked against the compiled auxiliary file by TXT.1--TXT.3.

\paragraph{Changes from the previous edition that affect claims.} (1) The value-of-observation bound \meqref{eq:voi} with its attaining family $\mathcal F_n$ (Proposition~\mref{prop:feedback}) and the attainment of the Sperner bound (Theorem~\mref{thm:antichain}) are new. (2) The exact threshold of a hold--probe--learn schedule, a parenthetical in the previous edition, is now the theorem (Theorem~\mref{thm:additive}), the additive law being its corollary under aligned worst losses, with an explicit misaligned instance (Example~\mref{ex:additive}) at which the additive formula overstates the threshold by one unit. (3) The contamination proposition carries the explicit two-stage counterexample to affine interpolation ($\tfr1{16}$ against $\tfr18$) and the convexity statement; the prior-contamination closed form (Proposition~\mref{prop:priorcontam}) is new. (4) The class scope of the timing obstruction, a remark before, is a corollary with proof (Corollary~\mref{cor:timing}), and the deficit profile on the certified grid is stated (Corollary~\mref{cor:grid}). (5) The two-floor instance and the asymmetric audit are examples with proofs; the noiseless-limit and delayed-instance-as-schedule remarks are new. (6) A paragraph cross-checking the minimax companion's interpolation record, which could not be regenerated from this paper's scripts, and a self-citation unrelated to the content have been removed. (7) The three tables of the previous main text are here (Tables~\ref{tab:grid}, \ref{tab:asym}, \ref{tab:deadline}) and are now generated, not typed; the five figures are replaced by four recomputed from the exact module.
\section{Symbol table}\label{si:symbols}

The table lists the symbols of the main text with the section in which each is introduced, and, where the object has a counterpart in the companion papers (the obstruction calculus \cite{Abaee2026calc} and the finite-certificate paper \cite{Abaee2026cert}), the reading used there. The companions are set-valued and deterministic: their objects are recovered here when every prior mass is positive on the compatible states and the kernels are deterministic (Theorem~\mref{thm:support}).

{\footnotesize
\begin{longtable}{@{}>{\raggedright\arraybackslash}p{0.21\textwidth}>{\raggedright\arraybackslash}p{0.49\textwidth}>{\raggedright\arraybackslash}p{0.24\textwidth}@{}}
\toprule
Symbol & Meaning (main-text section) & Companion reading\\
\midrule
\endfirsthead
\toprule
Symbol & Meaning (main-text section) & Companion reading\\
\midrule
\endhead
\bottomrule
\endlastfoot
$\calV$, $\bot$, $X^{\bot}$ & finite safe set, absorbing unsafe state, augmented state space (\S\mref{sec:setting}) & viability constraint set; exit\\
$A$, $Y$, $y_{\bot}$ & action set, observation set, exit announcement (\S\mref{sec:setting}) & controls; signals\\
$T$, $p(x,a)$ & transition kernel; one-step exit probability $T(\bot\mid x,a)$ (\S\mref{sec:setting}) & set-valued dynamics $F$ when deterministic\\
$g$, $h$ & observation law; deterministic observation map (\S\mref{sec:setting}) & observation map\\
$b$, $\supp b$, $b(\calV)$ & belief, its support, mass still alive (\S\mref{sec:setting}) & information state (= support)\\
$\tau(b,a,y)$, $\Prob(y\mid b,a)$ & Bayes update, outcome likelihood (\S\mref{sec:setting}) & post-state $\Post(B,a,y)$\\
$\Pi_{\seq}$, $\Pi_{\ol}(\ell)$, $\Pi_{\hold}(\ell)$, $\Pi_{\obs}$ & sequential, open-loop-on-window, hold-on-window and fully observed policy classes (\S\mref{sec:setting}) & regulators; hold-until-review regulator\\
$V^{\Pi}_k(b)$, $V_k(b)$ & safety value of class $\Pi$ at horizon $k$; sequential value (\S\mref{sec:setting}) & viable ($=1$) / nonviable ($<1$) verdict\\
$\Delta^{\Pi}_k(b)$ & deficit $1-V^{\Pi}_k(b)$ (\S\mref{sec:setting}, \S\mref{sec:deficit}) & ---\\
$\Gamma^{\mathrm s}_k(b)$, $S_k(b)$ & attaining set; value-one level of actions (\S\mref{sec:setting}) & $S_k(b)=\varnothing$ is an obstruction certificate\\
$\Gamma^{\Pi}_k$, $\alpha$ & $\alpha$-vector set and its elements, $V^{\Pi}_k(b)=\max_{\alpha}\langle\alpha,b\rangle$ (\S\mref{sec:setting}) & ---\\
$\calS^{\Pi}_k(b)$ & family of subsets of $\supp b$ kept alive by one policy of $\Pi$ for $k$ stages (\S\mref{sec:support}) & survivable sets\\
$\Wfb_k$, $\Wbl_k$ & feedback and blind survivable families (\S\mref{sec:support}) & $k$-step viable information states (feedback); common-action verdicts (blind)\\
$\mathcal F_n$ & $n$-state feedback family attaining the value-of-observation bound (\S\mref{sec:support}) & ---\\
$z_t$, $\theta$, $u_t$, $w$ & stock, hidden regime $\pm1$, control, prior mass of $\theta=+1$ in the delayed instance (\S\mref{sec:delayed}) & delayed-observation instance\\
$z_0$, $\Tobs$, $k$, $\ell$ & initial stock, review delay, horizon, effective window $\min(k,\Tobs)$ (\S\mref{sec:delayed}) & $\sigma^{*}=z_0-1$ exit time of the falling branch; timing certificate $z_0<1+\Tobs$\\
$V^{\hold}$, $V^{\ol}$, $V^{\seq}$, $V^{\obs}$ & class values of the delayed instance (\S\mref{sec:delayed}) & ---\\
viable / common-action / timing cell & the three cell classes of the certified grid (\S\mref{sec:delayed}) & verdict classes of the certified table\\
$\Tl$, $T$ & learning delay in the declining instance; generic delay in the schedule law (\S\mref{sec:deadline}) & ---\\
$d(u,\theta)$, $d_0$, $e_p$ & drift under action $u$ in regime $\theta$; worst hold loss; worst probe excursion (\S\mref{sec:deadline}) & ---\\
$t_p$, $z_p$, $\varepsilon$ & probe time, stock at the probe, probe error rate (\S\mref{sec:stochastic}) & ---\\
$p_{\mathrm{wrong}}(n)$ & probability that a majority of $n$ probe readings is wrong (\S\mref{sec:stochastic}) & ---\\
$\rho$, $\calD_\rho(p)$, $V^{\rho}_k$ & contamination level, $\varepsilon$-contamination ball of an outcome law, contaminated value (\S\mref{sec:dr}) & robust (adversarial) reading at $\rho=1$\\
$\ind{\cdot}$, $\tfr{a}{b}$ & indicator; exact rational constants (throughout) & ---\\
\end{longtable}}

\begin{thebibliography}{9}
\bibitem{Abaee2026calc} Abaee, A., 2026. An obstruction calculus for viability under incomplete observation. Figshare preprint, version 2. \url{https://doi.org/10.6084/m9.figshare.33716593.v2}.
\bibitem{Abaee2026cert} Abaee, A., 2026. Finite certificates of nonviability under incomplete observation: a continuous-to-finite bridge with exact witnesses. Companion manuscript, \url{https://github.com/MIKEAA2020/general-sustainability}.
\end{thebibliography}

\end{document}
